
In the traditional material  research, computational modelling had largely been applied for post-experimental rationalisation, which often starts from a very limited set of well defined chemical compositions and crystal structures. This was fundamentally limited by the high  costs associated with the computations, which are usually performed at the DFT-level of theory, and the computational power available at the time. With  more efficient implementations of the electronic structure theory codes, \emph{e.g.} on the Graphic Processing  Units (GPUs)\cite{hutchinson2012vasp,maintz2018strategies}, and increasing computational power on the  High Performance Computing (HPC) facilities, it has become increasingly popular and important to perform  samplings on the chemical and structural spaces for a specific material system to obtain its energy landscape as governed by some of the energetic merits discussed above, in order to facilitate the holistic determination and understanding of the materials' structures and properties, as well as proposing new material systems with unconventional properties. It also inspires the continuous development of cheaper yet accurate modelling methods, in particular, machine-learning models, for us to simulate materials in complex scales that were not accessible before. Here, we shall provide a brief account of the recent literature to showcase different areas of applications in which structural sampling in computational material studies is particularly important and useful.

\subsubsection{Structure Elucidations} Whilst XRD (X-Ray Diffraction) and TEM (Transmission Electron Microscopy) are the two mostly used and standard method for determining the structures of a given material, there exists scenarios when, a high quality crystalline sample cannot be obtained for detailed structural characterisations, or the material is in a form that only exists under extreme environments that cannot be characterised with the standard experimental techniques.

For examples, solid state nuclear magnetic resonance (NMR)\cite{hodgkinson2020nmr} has now become an important technique for determining the crystal structures of organic molecules that only exist in a power form either naturally or due to to the experimental limitations, which prevents crystals for diffraction being obtained. The NMR signals are sensitive to the local chemical environment around the NMR-active nuclei in the crystal, which is in term determined by the overall crystal packing. Therefore, the NMR signal can be considered as a \emph{fingerprint} of the crystal structure, where the later can be determined if this particular type of fingerprint can be decoded. By first applying organic crystal structure predictions\cite{price2014predicting} with cheap atom-atom potential to generate a  set of hypothetical crystal structures that are thermodynamically stable, and then subjecting them to periodic DFT calculations that enable the determinations of their corresponding (theoretical) NMR chemical shift signals,\cite{van1993first,ashbrook2016combining} one can reveal the crystal structure of the compound being the one with its theoretical NMR signal that best matches the experimental result.\cite{baias2013novo,selent2017clathrate} 

Similar approaches have also been adopted to facilitate the determination of the crystal structures for inorganic minerals and small gaseous molecules under thermodynamic conditions that are inaccessible under standard laboratory conditions, which is typically the case for geo- and planetary sciences.\cite{wang2014perspective} For instances, the crystal structures of the superionic phase of ice\cite{wang2011high} had been proposed under a high pressure ($>$14 Mbar) and zero temperature, which is believed to exist in the cores of planets such as Uranus and Naptune,\cite{guillot1999interiors} among the solid phases of many other small molecules.\cite{han2019predicting,cable2020properties,ennis2020mixed,maynard2022crystal} Using the complementary crystal structure prediction algorithms developed for inorganic materials,\cite{glass2006uspex} it has enabled geoscientists to gain more detailed structural insights into the phase behaviours of minerals that exists in the earth's mantle, such as the post-aragonite phase of \ce{CaCO3} at a pressure above 40 GPa,\cite{oganov2006high} the anisotropic post-perovskite structure of \ce{MgSiO3} at 120 GPa,\cite{oganov2004theoretical} as well as iron-carbides and iron-hydrides at 330-364 GPa.\cite{bazhanova2012fe} More recently, it has also been applied to provide structural models for understanding the mechanisms of high-temperature superconductivities in metal hydrides,\cite{salke2019synthesis,troyan2021anomalous,chen2021high} in which the anharmonic behaviour of lattice vibrations (see \cref{sec:anharmonic_vibs}) are also believed to be playing a critical contribution to the observed superconducting behaviours. 

\subsubsection{Phase Diagrams} Rather than being used for determining a particular structure of a material, more importantly, the enumeration of the global energy landscape for material with a specific composition is a critical steps towards constructing their phase diagram under different thermodynamic boundary conditions, which is important for assessing the (meta)-stabilities of different phases.\cite{nyman2015static,sun2016thermodynamic} These information may be further used to guide the discovery  of new materials,\cite{aykol2018thermodynamic} most importantly, to plan and design the pathways of solid state synthesis of complex multicomponent solids.\cite{hautier2012accuracy,mcdermott2021graph} 

For instance, the voltage that can be delivered by Li-battery may be correlated with the free-energies of  Li-intercalated compounds with different Li concentrations,\cite{van2013understanding} which is an important principle that can be applied in the computational designs of ion batteries, provided that one can map out the compositional dependent phase diagrams containing the thermodynamic stability information for all phases that may be formed by the target set of elements. \cite{doe2008first,albina2018finite} Elemental gallium possesses many intriguing physical properties, such as low melting point, density anomaly and the coexistence of covalency and metallicity. Recently, a full phase diagram has been constructed as facilitated by the neural network (NN) potential, showing excellent agreement with the experimental phase diagram, and the formation of the metastable $\beta$-Ga is s kinetically driven process over the formation of the thermodynamically stable $\alpha$-Ga.\cite{niu2020ab} Application of computational crystal structure prediction to the W-Se system\cite{zhu2020ab} reveals that the trigonal crystal structure (space group $P\bar{3}m1$) containing layers of van der Waals bonded \ce{WSe2} is the thermodynamic stable phase under ambient temperature and pressure. Increasing the pressures to 33 GPa, then subsequently 105 GPa, an orthorhombic $(Pmmn)$ and a cubic $(Pm\bar{3}m)$ phase can be stabilised with a notable change in their atomic coordination environments. 

\subsubsection{Phase Transitions and Lattice Dynamics}   Beyond the enumeration of different phases for a given material in its phase diagram, there is a strong desire to understand both the thermodynamic and atomistic origins behind the transitions between different phases for a given material, in order to gain better control in their phase behaviours and corresponding physical properties. For example, the inorganic halide perovskite \ce{CsPbI3}  has a suitable bandgap of 1.73 eV and has attracted immense intention to be developed as the next-generation solar cell, but it has a low phase stability, making it prone to transform into the wide-band-gap tetragonal phase at room temperature.\cite{swarnkar2016quantum} As such, computer simulation has become pivotal in providing detailed atomistic understanding behind the phase transition mechanisms in order to devise new strategies to stabilize the desired phase of technologically important materials under their standard operating condition. This generally requires the performance of long-time molecular dynamic (MD) simulations to achieve a sufficient sampling of the phase space in order to detect  the  `rare-events', such as the occurrence of a phase transition. These samplings could be particularly challenging when bond breakage and rearrangements, as well as anharmonic vibrations are involved,\cite{lahnsteiner2018finite,marronnier2017structural} making classical (harmonic) forcefield become incapable\cite{lahnsteiner2018finite} of producing the phase transition events, whereas DFT becomes too costly to achieve a thorough sampling of the relevant phase space. 

Such a challenge has gradually be overcome with the development of machine-learning potentials that is capable of reproducing the DFT-PES but at the cost of parameterised atom-atom forcefields. In the early years, these potentials are pre-trained specifically for a given compound (typically with a simple chemical composition) using high-dimensional NNs\cite{behler2011neural}, to examine, for example, the pressure-assisted  graphite-to-diamond phase transformation,\cite{khaliullin2011nucleation}  and the mechanical behaviour of the phase changing compound GeTe.\cite{sosso2013fast}.  More recently, with the development of better  kernel for atomistic machine learnings,\cite{bartok2013representing,de2016comparing} that is highly effective in interpolating the  PES,\cite{bartok2010gaussian} new scheme to train machine-learning potential on-the-fly from DFT-MD samplings have been developed. This method has since been successfully applied to study the entropy-driven temperature-dependent phase transformations in methylammonium lead iodide and its inorganic counterparts,\cite{jinnouchi2019phase} revealing systematically, that the cation radii are playing an important role in determining the phase transition temperatures of lead halide perovskites. The same method has also be applied to study the melting behavior of Si,\cite{jinnouchi2019fly} and \ce{ZrO2}.\cite{verdi2021thermal}

At a more fundamental level, the emergence of phase transition can be attributed to the presence of lattice instability that is characterised by the possession of imaginary phonon modes under the harmonic approximation,\cite{pallikara2022physical} reflecting the highly anharmonic nature of the lattice vibration. In particular, by continuously distorting the structure along the eigenvector of a given imaginary mode is another viable way to determine the phase transition pathway between a pair of crystal polymorphs.\cite{rahim2020polymorph,pallikara2021phase} Modelling the lattice vibrations, in particular, the anharmonic contributions to the lattice vibrations, also involves a sampling process, either in a brute-force enumeration or a Bayesian inference approach, to obtain a set of atomic displacements and their corresponding atomic forces, which are more elaborately discussed in \cref{sect:phonon_harmonic_and_anharmonic}. This often forms an initial setup to further sample the thermal induced fluctuations in the electronic structures of solids, which are important for characterizing the electron-phonon coupling behaviours that subsequently determine their properties such as the (photoexcited) carrier life-time and thermoelectric transports, particularly in mechanically soft materials.\cite{li2018spin,li2019anharmonicity,yang2020correlating,yang2022mapping,fratini2017map}  The modelling of electron-phonon couplings will be further discussed in \cref{sec:eph_coupling}.

\subsubsection{Intrinsic Defects and Chemical Dopants} Manipulating the physical properties of materials through controlling the defects and chemical dopants have been the central theme in traditional material science. As such, extensive efforts have been devoted to apply modelling as a key tool in facilitating the determinations and understanding of the atomistic structures and associated physical properties of defects and dopants in solids, which has been extensively discussed in many prior reviews.\cite{freysoldt2014first,anand2022convergent} Here, we focus on a few aspects for which the idea of sampling is particularly prevalent in the computational investigations of defects and dopants in solids.

Both point defects and dopants break the symmetry and long-range periodicity of a crystalline solids, and they are typically modelled with the `supercell' approach in DFT.\cite{makov1995periodic} For this, it is often necessary to minimize the artefacts induced by the unphysical Coulombic interactions among the periodic images of the defects/dopants by either sampling supercells with increasing sizes to extrapolate the asymptotic limit at infinite cell,\cite{kantorovich1999elimination} or applying correction to the results obtained from small supercell.\cite{freysoldt2009fully} The increase in the computational power has propelled the community to investigate the defect and dopant chemistry in materials with more complex chemical compositions and crystal structures, meaning defects/dopants across a wider range of the periodic table, as well as multiple different lattice occupancies, can be sampled in the host materials. The latter is important to investigate the possible configurational entropic stabilisation of defects/dopants in the host lattice, which counterbalances the enthalpic penalties in their formations. These need to be further accompanied with, as shown in \cref{eq:formation_defect_chem_pot}, the investigations of charge-dependent stabilities of defects/dopants, which consequently determines their electronic energy levels in the band structures. 

Some examples that best illustrate the above aspects includes: The modelling the defect chemistries in lithium  solid electrolytes (\ce{Li7La3Zr2O12}) with a complex garnet-type crystal structure,\cite{miara2015first,squires2019native} cobaltate thermoelectrics with incommensurate layered structures,\cite{rebola2012first,lei2014synthesis} as well as complex oxides (\emph{e.g.} \ce{SrCoO}$_{x}$\cite{jeen2013topotactic} and \ce{SrFeO}$_{x}$\cite{nguyen2020topotactic}) that can undergo  topotatic phase transformations which can be tuned by the oxygen vacancy concentration. In the last case, it implies that the defect formation energy for a given compound at a given concentration needs to be additionally sampled across a few different phases in order to identify the thermodynamically most likely phase to accommodate oxygen vacancies with a specific concentration. The occurrence of defect complexes in solids, particularly hydrogens,\cite{morris2008hydrogen,lange2010native}  further combinatorially increases the size of the search space to identify the lowest-energy lattice configurations for these complexes. In this regards, whilst most of the prior computational studies constructed supercells containing defects/dopants through ion replacements on lattice site(s) followed by an energy minimisation to the \emph{local} minimum, a recent investigation\cite{mosquera2022identifying} has highlighted the fact that this approach does not guarantee the reach of their energetic ground-state defect/dopant configurations, rendering further configuration samplings around the local minimum should be further incorporated into the protocol for modelling defects/dopants in crystalline solids.


The above discussions have ignored the temperature contributions to the thermodynamic and physical properties of dopants or defects, which may come from a few different aspects that will be briefly discussed here. Firstly, the equilibrium concentration $n_{D}^{q}$ for a given dopant/defect $D$ in the charge state $q$ is given by minimizing the free energy of defect formation with respect to $n_{D}^{q}$,
\begin{equation*}
    G(T,p)=\sum_{q}n_D^q \Delta G_D^{q}(T,p)-G^{0}(T,p)-T\Delta S_{conf},
\end{equation*}
in which $G^{0}(T,p)$ and $\Delta G_D^{q}(T,p)$ are the temperature-dependent Gibbs formation energies of solid without and with the charged dopant/defect that may or may not further include the vibrational  (see \cref{sect:phonon}) and electronic contributions to the system entropies. Note that the formation of point defects in solids is expected to change the bonding strengths in a solid, and hence leads to a change in its phonon eigenspectrum and the vibrational entropies determined therefrom. The reduction in the crystal symmetry due to the presence of point defects also increases the number of symmetry inequivalent atomic displacements that need to be sampled for obtaining the phonon eigenspectrum.

Nevertheless, assuming all the defects are non-interacting hence different defect configurations are energetically degenerate, then the configuration entropy $S_{conf}$ is only dependent on the available dopant/defect sites ($N_D$) in the crystal and their corresponding occupancy $(n_{D}^{q})$.  (This is generally not the case, see the following discussions on high-entropy alloys) In this case, the minimisation problem $\partial G(T,p)/\partial n_{D}^q=0$ can be solved analytically, leading to the following Boltzmann-like expression for the equilibrium point defect/dopant concentration:\cite{kasamatsu2011theoretical}
\begin{equation*}\label{eq:defect_concentration}
    \left[D^q\right]=\frac{N_{D}}{V}\frac{\exp[-\beta\Delta E_f(D^{q})]}{1+\sum_{q^\prime} \exp[-\beta\Delta E_f(D^{q^\prime})] },
\end{equation*}
where $V$ is the cell volume per formula unit, $\beta=1/k_{B}T$ with $k_{B}$ being the Boltzmann constant, $q^{\prime}$ denotes all possible charge state for the defect/dopant, and $\Delta E_{f}$ being its corresponding formation energy [see \cref{eq:formation_defect_chem_pot}]. 

Secondly, in (strongly) anharmonic solids with notable electron-phonon coupling, additional molecular dynamic simulations may also be needed in order to model the thermally induced fluctuations in the defect/dopant electrons energy levels, in order to gain a more thorough understanding on the materials' physical properties close to the normal operating conditions, which may deviate significantly from the results obtained from static calculations.\cite{cohen2019breakdown,li2019influence,qiao2020extending,chu2020low}

\subsubsection{From Binary to High-Entropy Alloys}\label{sect:HEA}
Whilst defects and dopants are important in material designs, they are often present in the host lattice in low concentrations that do not fundamentally change the physical properties of the host materials. By \emph{randomly} intermixing two more more principal elements (often with similar chemical/physical properties) in high or equimolar concentrations to form homogeneous solid solutions (or alloys), it offers material scientists a  strategy to obtain new, high-entropy materials\cite{george2019high,oses2020high} with unusual chemical bonding states between pairs of elements, which otherwise may not be  stable as a binary compound. In these type of materials, the atoms for the same element occupying different lattice sites are surrounded by atoms of a different set of elements, which `confuses' the atoms to find an enthalpically favored pathway to `phase-segregate' into multiple different phases made from subsets of the constituent atoms. 

Thermodynamically, the stability of a multicomponent alloy is determined by its temperature-dependent formation free energy $\Delta G_f(T)=\Delta H_f-T\Delta S_{conf}$. Due to the presence of unusual bonding state, generally we have $\Delta H_f>0$, implying the system can only be stabilised if the configuration entropy that is contributed by different energetically (similar) atomic configurations in a multicomponent alloy, $T\Delta S_{conf}$ is larger than $\Delta H_f$. As pointed out in the seminal work by \citet{yeh2004nanostructured} and the discussions above on defects/dopants, when intermixing $n$ different principal elements \textit{in equimolar ratio} that do not have mutual interactions amongst each other, $S_{conf}=R\ln n$ (with $R$ being the ideal gas constant), showing $TS_{conf}>\Delta H_f$ is a condition that may be met when $n$ is large, such that a high-entropy system is often regarded as those with $n\geq 5$, although  $S_{conf}$ is not non-zero even in a binary mixture. However, in real materials, the principal elements interacts with each others, such that every unique atomic configuration of the principal elements in the lattice may have a different total energy $E^{(i)}$. This must be taken into consideration when determining the configuration entropy in a multicomponent alloy, which can be obtained through the Boltzmann relationship $G_{conf}=-k_{B}T\ln Z_{conf}$, where the configurational partition function
\begin{equation}\label{eq:configuration_partiton_function}
\mathcal{Z}_{conf}=\sum_{i=1}^{\mathcal{N}}g_{i}\exp[-\beta E^{(i)}],
\end{equation}
in which $g_i$ is the degeneracy factor. If one wishes to accurately determine the configurational entropy, then
\begin{align}\label{eq:conf_entropy}
    S&_{conf}=-\left(\frac{\partial G_{conf}}{\partial T}\right)_V\nonumber\\
    &=\frac{\partial}{\partial T}(k_{B}T\ln Z_{conf})_V\nonumber\\
    &=k_B \ln \sum_{i=1}^{\mathcal{N}}\exp[-\beta E^{(i)}]+\frac{1}{T}\frac{\sum_{i=1}^{\mathcal{N}}E^{(i)} \exp[-\beta E^{(i)}]}{\sum_{i=1}^{\mathcal{N}}\exp[-\beta E^{(i)}]},
\end{align}
assuming $g_i=1$. Strictly speaking, at finite temperature, the energy per configuration $E^{(i)}$ which enters \cref{eq:conf_entropy} should really be the Helmholtz free energy [\cref{eq:helmholtz}] that takes into account the vibrational entropic contributions at finite temperature.\cite{van2002effect} However, computing the Helmholtz free energy individually for each configuration across a wide set of configurations for high entropy materials with a brute-force approach is generally not feasible. Some recent attempt\cite{yang2025accounting} along this line can be seen with the use of deep-learning potentials that provide significant acceleration in atomistic simulations for new materials.  
Before proceeding, we note that the configurational entropy from the rotational degree-of-freedom is also important contribution to the free energies of flexible molecules, the determination of which has been a long time focus in theoretical chemistry.\cite{karplus1981method,hnizdo2008efficient,pracht2021calculation}
 
In solid-state materials, obtaining a converged estimation of $Z_{conf}$ based on \cref{eq:configuration_partiton_function} is a genuine sampling challenging, particularly when DFT is used for determining $E^{(i)}$. Firstly,  large supercell containing $N$ sites that can be replaced by up to $K$ different elements is needed to accurately simulate \emph{each} atomistic model of a high-entropy alloy, which itself is computationally demanding. Note that although there exists method, such as the virtual crystal approximation (VCA),\cite{bellaiche2000virtual} that enables the simulation of alloys using small unit cells using the compositionally averaged crystal potential, it will not allows us to explicitly enumerate individual $E^{(i)}$ term that goes into \cref{eq:configuration_partiton_function} for the corresponding determination of $Z_{conf}$. Secondly, the total number of configurations $\mathcal{N}=N!/\sum_{j=1}^k N_j!$,\cite{he2017formation} where $N_j$ is the number of sites that can be occupied by element $j$ and $\sum_{j=1}^k N_j=N$, showing the number of symmetrically independent configurations that need to be sampled grows factorially as $N$ (or equivalently, the supercell size) increases. Finally, the number of possible combinations of elements that can be obtained by choosing them from a pool of total $m$ elements, forming alloys with $i=2,3,\dots, m$ principal elements is $\sum_{i=2}^{m} C_{i}^{m}$, where $C_{i}^{m}$ is the number of combinations of $i$ items from a total of $m$. While this means we have a large chemical space (e.g. 8165 combinations in total for $m=13$) to search for new alloys, it also imposes additional challenges in computational material discoveries.

The enormously large number of configurations that can exist in a high-entropy alloy means there is a large variety of the local (metastable) chemical bonding environments that may not be found in the conventional materials, which may lead to outstanding catalytic properties. It also means that the interstitial crystal sites will exhibit a distribution in their sizes, with each being surrounded by a different set of ions, providing a higher probability for the high-entropy alloy to have the interstitial sites that are just big enough to accommodate hydride ions, making them to be ideal candidates to be further developed as hydrogen storage materials.\cite{marques2021review} In this regard, that is additional to determining the degree of entropy stabilisation in high-entropy alloys, sampling is important  in the computational design of functional high-entropy alloys as their overall properties, being either the catalytic performances or the hydrogen storage capacities, will be a statistical average over the probabilities that different atomic coordination environments or the interstitial void spaces will occur in their crystal lattice. For instances, by combining DFT and ML model, \citet{pedersen2020high} sampled the adsorption energies of CO on the surfaces of two high-entropy alloys \ce{CoCuGaNiZn} and \ce{AgAuCuPdPt} over 1000+ different local chemical environments (as characterised by the neighbouring atoms) of the adsorption sites, demonstrating a broad distribution in the absorption energies from weak non-binding ($\sim 0.5$ eV) to strong binding ($\sim -2$ eV) could be observed in a single high-entropy alloy. Similar approach has been applied to the design of high-entropy alloy catalysts for the oxygen reduction reactions (ORR),\cite{batchelor2021complex,pedersen2021bayesian} \ce{NH3} decomposition\cite{xie2019highly,yao2020computationally} and  \ce{CO2} reduction.\cite{chen2019identifying} Lastly, both the atomic binding energies and averaged atomic volumes (which can be correlated with the volumes of the interstitial sites) have been incorporated in the screening approach recently proposed by \citet{witman2021data} to optimize the compositions of the high-entropy alloys for hydrogen storage.

Inspired by the concept of high-entropy metal alloys, there is an increasing focus in the design of entropy stabilised functional ceramics, that are facilitated by the high-throughput computational screening. For instances, the emergence of methylammonium lead iodide (MAPbI$_3$) photoabosrbers\cite{zhou2014interface} has stimulated a significant body of works being devoted to the theoretical investigations of mixing different cationic/anionic species in the host halide perovskites to achieve tunable photovoltaic properties, including band gaps and carrier lifetimes, in a configurational-entropy-stabilised random alloy  that is resistant to phase separations.\cite{brivio2016thermodynamic} Some  examples include: A-site cation mixed perovskite (MA$_{x}$FA$_{1-x}$)PbI$_{3}$ (FA=formamidinium),\cite{senno2021mixed,schelhas2019insights} B-site cation mixed perovskite Cs(Pb$_{x}$Sn$_{1-x}$)X$_{3}$ (X=Cl, Br, I),\cite{yang2020mapping} FA(Pb$_{x}$Sn$_{1-x}$)I$_{3}$,\cite{kahmann2022compositional} X-site anion mixed perovskites CsPb(I$_{x}$Br$_{1-x}$)$_3$,\cite{yang2020correlating}  {C}{H}$_3${N}{H}$_3${P}b({I}$_{1-x}${B}r$_{x}$)$_3$,\cite{brivio2016thermodynamic} hetero-valent-cation-doped CsPbIBr$_2$,\cite{hu2022valence}  as well as low dimensional antimony perovskites Cs$_3$Sb$_2$X$_{y}$I$_{9-y}$.\cite{yang2021composition}  Anion mixing has also been applied to enhance the photovoltaic properties of kesterites, such as  Cu$_2$ZnSn(S,Se)$_4$.\cite{shang2014cation}
Because of the significant electronic interactions amongst two types of different but chemically similar ions occupying the same lattice site, resulting in a non-linear behaviour in their thermodynamic stabilities and electronic properties that deviates from the linear Vergard's law for a binary mixture,\cite{dalpian2019formation} performing configurational samplings in the computational studies of these perovskite alloys thus become particularly important and interesting in understanding their fundamental chemical and physical properties.  Alloying has also been adapted as a strategy to tune the electronic and phononic properties of layered two-dimensional materials such as {S}n({S}$_{1-x}${S}e$_x$)\cite{gunn2019thermodynamics,skelton2020lattice}  MXenes,\cite{han2020tailoring,nemani2021high} and transition metal dichalcogenides\cite{cavin20212d} for optical, thermoelectric and catalytic applications. The charge states of the transition metal B-site cations in complex oxide perovskites are sensitive to the oxygen stoichiometries, for which the subtle interplays between these two factors can be used to design thermal storage materials by interconverting between heat and the chemical energies in the B-O bonds,\cite{lei2022redox} This bonding strengths can be further manipulated by changing the B-site coordination environment through A-site alloying.\cite{sai2020exploring} 

It is important to note that the configurational entropy reflects the topology of the PES that is `accessible' at a given temperature $T$. Due to the high cost for computing the configurational entropy,  the term $E^{(i)}$ that enters \cref{eq:conf_entropy} is usually the enthalpy of the configuration corresponding to the local energy minimum at 0 K. However, at finite $T$, a configuration is not frozen at its minimum energy  atomic arrangement taking a single energy state of $E^{(i)}$, but constantly vibrating around the minimum within the local energy basin. In other words, there is no reason why the phononic contribution (\cref{sect:phonon}) to the stability of individual configuration should not be accounted for when calculating the configurational entropy. As formulated by \citet{van2002effect}, this can be achieved by replacing $E^{(i)}$ by its corresponding Helmholtz free energy [\cref{eq:helmholtz}, which is effectively the expected energy for a  given configuration that is confined and vibrating in its energy basin at $T$] in \cref{eq:conf_entropy} to effectively include the vibrational entropy $S_{vib}$ in the configurational entropy. Such a calculation would be prohibitively expensive with DFT, and  solution to such a challenge using machine-learning atom-atom potential has emerged recently, \cite{tolborg2023low} which, for the solid solution of Na$_{1-x}$K$_{x}$Cl, including $S_{vib}$ has resulted in the theoretically calculated miscibility dome for this particular solid solution become more closely resembling the one that was determined experimentally.

\subsubsection{Amorphous Solids} The amorphous (glassy) state of solids constitutes a higher level up in the structural complexities of solids. Compared to the high-entropy alloys, not only is the long-range periodicity absent in the amorphous solids, the short-range atomic environments as characterised by the bond lengths and angles, particularly, the coordination numbers, can be completely different even between two neighbouring atoms. These diverse chemical bonding environments have contributed to unique strengths, corrosion resistance and chemical activeness that benefit the application of amorphous materials as coating layers,\cite{joshi2016amorphous} catalysts,\cite{anantharaj2020amorphous}  photovoltaics,\cite{stuckelberger2017progress} and many others. A detailed understanding of the atomic bonding environments in amorphous solids will inevitably play a critical role to facilitate the understanding of the origin behind their unique physical and chemical properties. However, because there exists limited (experimental) knowledge of the detailed atomistic structures for amorphous solids, some types of unbiased sampling will be necessary for us to obtain an ensemble of realistic atomistic models of the amorphous solids. 

Computationally, this can be achieved from running molecular dynamic (MD) simulations following the typical temperature profile of how amorphous solids are formed experimentally, that is superheating the material and equilibrating it in its liquid state, followed by a rapid quenching to freeze the material into a snapshot of its nonequilibrium structure. A few challenges exist in performing this type of samplings. Similar to the high-entropy alloys, large supercells containing $10^3$-$10^5$ atoms are needed for simulating amorphous solids that are lack of the long-range periodicity. Although systems of this size can typically be handled with classical forcefields, they are typically parameterised (biased) against a subset of prototypical structures, resulting them being incapable of accurately reproducing the diverse chemical bonding environments that are revealed from detailed experimental characterisations.\cite{de2019transferability} Whilst this parameterisation problem can be mitigated by running the simulation with DFT, the poor scaling behaviour of DFT has significantly limited the system sizes (up to a few hundred atoms) that can be modelled, which must be further accompanied by a fast quenching rate ($10^{14}$ K/s).\cite{cooper2000density} On the other hand, a cooling rate that is three orders of magnitudes slower is required, for example, to correctly reproduce the bonding environments in amorphous silicon\cite{deringer2018realistic} as revealed by the solid-state NMR measurements. 

Significant progress in overcoming the aforementioned challenges to obtain realistic models of amorphous solids comes with the developments of ML forcefields that are trained on the DFT-determined energies and forces for a particular material with a diverse range of coordination environments, particularly the works on the application of the GAP (Gaussian Approximation Potential) model\cite{bartok2010gaussian} to amorphous systems as led by Cs\'{a}nyi and Deringer. For example, in one of the earliest works,\cite{deringer2017machine} GAP is applied to simulate the structures of both the bulk amorphous carbon (a-C) and its surfaces. Unlike the forcefield model, GAP was successful in reproducing the experimentally determined contents of $sp^3$-coordinated C atoms across the spectrum from low (1.5 gcm$^{-3}$) to high (3.5 gcm$^{-3}$) density a-C. It also shows that the extent of surface graphitisation increases with the annealing temperature, that is also in compliance with the experimental observations. The robustness of the model was subsequently put to test on simulating the depositions of a-C films with incident ion energies from 5 to 100 eV.\cite{caro2018growth,caro2020machine} The simulations showed that low impact energies give rise to low density a-C that is dominated by graphite-like $sp^2$-C, whereas high impact energies lead to the growth of high density a-C with predominately diamond-like $sp^3$-C. The simulation clearly revealed that, contradicting to the proposed `subplantation' mode, the high-density a-C is grown via a `peening' mode where the high-energy carbon impacted onto the surface generates a pressure wave that induces densifications away from the impact site. Using a cut-off radius (1.9 $\mbox{\AA}$) based criteria, it was also found that a small but non-negligible number of C atoms was in a fivefold coordination environment resulting from high energy deposition. Detailed bonding analysis from quantum mechanical calculation reveals that the fifth C atom, that is the atom furthest from the center atom but residing within the cut-off radius, shows negligible bonding strength with the center C. This means that the fivefold coordinated C is an artifact from how the coordination sphere was algorithmically identified, rather than a defect of the GAP model. 

Subsequently, the GAP model is  applied to investigate the amorphous silicon (a-Si) system using supercell containing $10^5$ atoms.\cite{deringer2021origins} The  vitrification process of Si is simulated from a liquid state at 1,500 K to a lower temperature of 1,175 K over a 1 ns-long MD trajectory, which clearly reveals the gradual decrease in the Si coordination number from $>8$ to the normal 4-fold tetrahedral coordination, that is accompanied with the lowering of the atomic energy, indicating the stabilisation effect. For the first time, this work also provided detailed atomistic insight into the pressure induced structural transitions in a-Si, showing that, at 13 GPa, the low-coordination a-Si transforms into a transient VHDA (very high density amorphous) state with coordination number $\geq 10$. Further increasing the pressure to 15 GPa, crystallisation is initiated with the emergence of crystalline domains that are primarily made of 6-fold coordinated Si, where the system fully transforms into a polycrystalline state at 20 GPa. More uniquely, the pressure-induced crystallisation is accompanied with an electronic  transition to the  metallic phase at pressure $\geq 12$ GPa, which is revealed by applying an ML model that enables the prediction of the electronic density-of-states (DOS) over an ensemble of disordered structures.\cite{mahmoud2020learning}

\subsubsection{Surfaces, Interfaces and Heterogeneous Catalysts} The creation of a surface breaks both the translational symmetry, as well as the equilibrium atomic coordination environments in the bulk crystal structure. Consequentially, this brings the system to a higher energy state, and the magnitude of this energy increase determines the thermodynamic stabilities of surfaces, which can be affected by factors such as the surface orientations and termination chemistries or stoichiometries, meaning that these factors need to be thoroughly sampled in order to obtain a comprehensive surface stability diagram for a given material [\emph{c.f.} \cref{fig:energy_landscape_examples}(c-d)]. To minimize such an energetic penalty in the surface creations, the topmost atomic layer will typically undergo surface reconstructions \emph{in vacuo} to compensate the surface dangling bonds, and some of these reconstructions can lead to a large surface unit cell\cite{kruger2006generalized,battaglia2009structure} that imposes additional challenges in their computational investigations. By joining two surfaces, from either different materials to form a heterointerface, or the same material to form grain boundaries, new physical properties may emerge. Although there exist a large body of theoretical works being devoted to  epitaxial heterointerfaces, less is so for the studies of grain boundaries, as the (often unknown) atomistic structures at the grain boundary regions are expected to be significantly different from the bulk, that is imposed by the crystallographic orientations and the grain boundary angles.\cite{chua2010genetic} This requires significant computational power to determine the lowest energy grain boundary structure through sampling, especially when the supercell size for sampling is generally inversely proportional to the grain boundary angles. Finally, accurate surface models are the pre-requisite for theoretically  investigating the  catalytic properties of heterogeneous catalysts, in which samplings are further needed to determine the catalytic active sites on the surfaces, the possible number of choices for which increases with lowering surface symmetry and degree of chemical complexities (such as the high-entropy alloys discussed above). 



We note that the sampling can become more challenging if the surface adsorbates contain additional internal degrees of conformational flexibility and different charge states, given that the simulation of hybrid organic/inorganic interface itself is already considered to be theoretically challenging.\cite{draxl2014organic} For example, as demonstrated recently by \citet{maksimov2021conformational}, in which they have sampled more than 1000 configurations of arginine on three different metal surfaces, from which the results revealed the accessible molecular surface area of surface adsorbed arginine is strongly depending on the protonation state of the arginine molecule. 

\subsubsection{Closing Remarks}  
The above examples show that sampling is an overarching theme in modern theoretical material science, which is enabled by the increasing computational power and the development of ML methodologies. As a result, rather than focusing on examining a specific physical property $\mathcal{A}$ for material with a given atomic configuration $\bm{R}$, as how computation was used to be applied in supplementing experimental observations, there is an increasing capability for material theoreticians to better predict the properties of a specific material system by taking an `ensemble averaged' to determine the materials' expected property $\langle \mathcal{A}\rangle_{T}$ at their working temperature $T$ as
\begin{equation}\label{eq:general_sampling_1}
\langle \mathcal{A}\rangle_{T} = \int_{\bm{R}} \mathcal{A}(\bm{R})\mathcal{P}\left[E(\bm{R}),T\right]\mbox{d}\bm{R} \quad \mbox{(continuous)},
\end{equation}
or
\begin{equation}\label{eq:general_sampling_2}
\langle \mathcal{A}\rangle_{T} = 
\sum_i \mathcal{A}(\bm{R}_i)\mathcal{P}\left[E(\bm{R}_i),T\right] \quad \mbox{(discrete)}, 
\end{equation}
in which $\mathcal{P}$ denotes the corresponding statistical distribution of atomic configuration $\bm{R}$ with an energy of $E(\bm{R})$. It is obvious that the determination of $\langle \mathcal{A}\rangle_{T}$  is nowadays the key challenge that one wishes to tackle in modern computational material science, which heavily relies on our modelling capabilities to: (1) generate a database of structures $\left \{\bm{R}\right\}$ for new materials, (2) evaluate their corresponding stabilities $E(\bm{R})$ and (3) the associated properties $\mathcal{A}(\bm{R})$. The following three sections of this review will therefore be devoted for discussing these three aspects in more detail.
